\vspace{-2mm}
\subsection{LLM's Diverse Capabilities in Video Generation}
\label{sec:capabilities_whole}
% \vspace{-2mm}

This subsection presents several capabilities we discover from the pretrained \modelname{}, shedding light on the LLM's promising potential in video generation. By combining the flexibility of our autoregressive language model to perform diverse tasks such as extending video in time, inpainting, outpainting, and stylization, \modelname{} accomplishes multiple tasks in a unified model.


\vspace{-4mm}
\paragraph{Coherent long video generation and image-to-video.}

\begin{figure}[t]
\addtocounter{figure}{-1}

\begin{subfigure}{1.0\linewidth}
\centering
\videoframe{figures/video_frames/astronaut_00001.jpg}
\videoframe{figures/video_frames/astronaut_00029.jpg}
\videoframe{figures/video_frames/astronaut_00060.jpg}
\videoframe{figures/video_frames/astronaut_00081.jpg}

\end{subfigure}

\captionof{figure}{\textbf{10-Second long video generation example.} By predicting 1-second video segments from an initial 1-second clip, \modelname{} can iteratively generate videos of extended lengths.}
\label{fig:long_video}
\vspace{-4mm}
\end{figure}

\begin{figure}[t]
\addtocounter{figure}{-1}

% \begin{subfigure}{1.0\linewidth}
% \centering
% \videoframe{figures/video_frames/milk_drop_00001.jpg}
% \videoframe{figures/video_frames/milk_drop_00004.jpg}
% \videoframe{figures/video_frames/milk_drop_00008.jpg}
% \videoframe{figures/video_frames/milk_drop_00014.jpg}
% {\textbf{Animated from Photograph}
% \ \\
% \ \\}
% \end{subfigure}

\begin{subfigure}{1.0\linewidth}
\centering
\videoframe{figures/video_frames/flag_00001.jpg}
\videoframe{figures/video_frames/flag_00006.jpg}
\videoframe{figures/video_frames/flag_00011.jpg}
\videoframe{figures/video_frames/flag_00017.jpg}
{\textbf{Animated from historical photo}
\ \\
\ \\}
\end{subfigure}

% \vspace{-3mm}
\begin{subfigure}{1.0\linewidth}
\centering
\videoframe{figures/video_frames/ship_00001.jpg}
\videoframe{figures/video_frames/ship_00003.jpg}
\videoframe{figures/video_frames/ship_00009.jpg}
\videoframe{figures/video_frames/ship_00015.jpg}
{\textbf{Animated from painting}
\ \\}
\end{subfigure}
% \vspace{-3mm}
\captionof{figure}{\textbf{Examples of videos animated from still images plus text prompts tailored to each initial image.}}
\label{fig:animated_images}
\end{figure}

A benefit of an decoder-based language model is that it pairs well with autoregressively extending generation in time. 
%
We present two different variants: Generating longer videos and converting images to videos. 
Encoding the first frame independently allows us to convert any image into the initial frame of a video without padding. Subsequent frames are generated by predicting remaining tokens, transforming the image into a video as shown in \Cref{fig:animated_images}\footnote{For image-to-video examples we source images from Wikimedia Commons: https://commons.wikimedia.org/wiki/Main\_Page}.

This results in the capability to generate videos longer than 10 seconds or to allow users to iteratively extend video clips based on previously generated video, and produces temporally consistent videos without significant distortion.
Such capabilities are rarely observed in contemporary diffusion models.

\vspace{-4mm}
\paragraph{Zero-shot video editing and task chaining.} \label{sec:capabilities}
% \begin{figure}[t]
\addtocounter{figure}{-1}

\begin{subfigure}{1.0\linewidth}
\centering
\videoframe{figures/video_frames/goat_original_00001.jpg}
\videoframe{figures/video_frames/goat_original_00006.jpg}
\videoframe{figures/video_frames/goat_original_00011.jpg}
\videoframe{figures/video_frames/goat_original_00017.jpg}
{\textbf{Original Video}
\ \\
\ \\}
\end{subfigure}


\begin{subfigure}{1.0\linewidth}
\centering
\videoframe{figures/video_frames/goat_mask_00001.jpg}
\videoframe{figures/video_frames/goat_mask_00006.jpg}
\videoframe{figures/video_frames/goat_mask_00011.jpg}
\videoframe{figures/video_frames/goat_mask_00017.jpg}
{\textbf{Masked Video}
\ \\
\ \\}
\end{subfigure}


\begin{subfigure}{1.0\linewidth}
\centering
\videoframe{figures/video_frames/goat_dragon_00001.jpg}
\videoframe{figures/video_frames/goat_dragon_00006.jpg}
\videoframe{figures/video_frames/goat_dragon_00011.jpg}
\videoframe{figures/video_frames/goat_dragon_00017.jpg}
{\textbf{Inpainted Video from Text Prompt} \\
\textbf{Prompt:} A blue dragon walking along a ridge
\ \\}
\end{subfigure}

\captionof{figure}{
\textbf{\textbf{Example of a task  the model was not explicitly trained for} -- video editing via text conditioned inpainting.}
}
\label{fig:inpainting_dragon_goat}
\end{figure}
\begin{figure}[tp]
\addtocounter{figure}{-1}
\centering
\begin{subfigure}{1.0\linewidth}
\centering
\videoframe{figures/video_frames/mona_yawn_lores_00001.jpg}
\videoframe{figures/video_frames/mona_yawn_lores_00006.jpg}
\videoframe{figures/video_frames/mona_yawn_lores_00011.jpg}
\videoframe{figures/video_frames/mona_yawn_lores_00017.jpg}
{\textbf{Animated from still image}
\ \\
\ \\}
\end{subfigure}
% \vspace{-4mm}


\begin{subfigure}{1.0\linewidth}
\centering
\videoframe{figures/video_frames/mona_yawn_snow_00001.jpg}
\videoframe{figures/video_frames/mona_yawn_snow_00006.jpg}
\videoframe{figures/video_frames/mona_yawn_snow_00011.jpg}
\videoframe{figures/video_frames/mona_yawn_snow_00017.jpg}
{\textbf{Stylized video} \\
\textbf{Prompt:} An oil painting of a snowman with a red hat opening their mouth to yawn
\ \\}
\end{subfigure}
\captionof{figure}{\textbf{Example of zero-shot video editing via task chaining} (text conditioned image-to-video and stylization) 
%-- the original painting is first animated via a text prompt and then stylized via another text prompt.
}
\label{fig:animate_stylize}
\vspace{-2mm}
\end{figure}


% %\begin{figure*}[h]
% \begin{figure}[t]
% % \addtocounter{figure}{-1} %Prevent subfigures increasing main figure count  % for some reason, this needed to be commented out to get this figure to have a different number than the preceeding figure.
% \centering
% % Arc shot
% \begin{subfigure}[b]{1.0\linewidth}
% %{0.5\textwidth}
% \centering
% \videoframe{figures/video_frames/train_original_00001.jpg}
% \videoframe{figures/video_frames/train_original_00006.jpg}
% \videoframe{figures/video_frames/train_original_00011.jpg}
% \videoframe{figures/video_frames/train_original_00017.jpg}
% %\captionof{figure}{Example of directed camera movement }
% %\end{figure}
% {\textbf{Original Video}
% \ \\
% \ \\}
% \end{subfigure}


% \begin{subfigure}[b]{1.0\linewidth}
% \centering
% \videoframe{figures/video_frames/train_outpaint_00001.jpg}
% \videoframe{figures/video_frames/train_outpaint_00006.jpg}
% \videoframe{figures/video_frames/train_outpaint_00011.jpg}
% \videoframe{figures/video_frames/train_outpaint_00017.jpg}
% {\textbf{Outpainted Video}
% \ \\
% \ \\}
% \end{subfigure}


% \begin{subfigure}[b]{1.0\linewidth}
% \centering
% \videoframe{figures/video_frames/train_food_00001.jpg}
% \videoframe{figures/video_frames/train_food_00006.jpg}
% \videoframe{figures/video_frames/train_food_00011.jpg}
% \videoframe{figures/video_frames/train_food_00017.jpg}
% {\textbf{Stylized Video} \\
% \textbf{Prompt:} A gingerbread and candy train on a track
% \ \\}
% \end{subfigure}

% \caption{\textbf{Example of zero-shot video editing via task chaining (outpainting and stylization)} -- the original video is first outpainted and then stylized via a text prompt.\lu{We can see a clear artifact in the middle row over the train (white block). This image may raise question.}}
% \label{fig:outpaint_stylize}
% \end{figure}

% %\begin{figure*}[h]
\begin{figure}[t]
% \addtocounter{figure}{-1} %Prevent subfigures increasing main figure count  % for some reason, this needed to be commented out to get this figure to have a different number than the preceeding figure.
\centering
% Arc shot
\begin{subfigure}[b]{1.0\linewidth}
%{0.5\textwidth}
\centering
\videoframe{figures/video_frames/train_original_00001.jpg}
\videoframe{figures/video_frames/train_original_00006.jpg}
\videoframe{figures/video_frames/train_original_00011.jpg}
\videoframe{figures/video_frames/train_original_00017.jpg}
%\captionof{figure}{Example of directed camera movement }
%\end{figure}
{\textbf{Original Video}
\ \\
\ \\}
\end{subfigure}


\begin{subfigure}[b]{1.0\linewidth}
\centering
\videoframe{figures/video_frames/train_outpaint_00001.jpg}
\videoframe{figures/video_frames/train_outpaint_00006.jpg}
\videoframe{figures/video_frames/train_outpaint_00011.jpg}
\videoframe{figures/video_frames/train_outpaint_00017.jpg}
{\textbf{Outpainted Video}
\ \\
\ \\}
\end{subfigure}


\begin{subfigure}[b]{1.0\linewidth}
\centering
\videoframe{figures/video_frames/train_food_00001.jpg}
\videoframe{figures/video_frames/train_food_00006.jpg}
\videoframe{figures/video_frames/train_food_00011.jpg}
\videoframe{figures/video_frames/train_food_00017.jpg}
{\textbf{Stylized Video} \\
\textbf{Prompt:} A gingerbread and candy train on a track
\ \\}
\end{subfigure}

\caption{\textbf{Example of zero-shot video editing via task chaining (outpainting and stylization)} -- the original video is first outpainted and then stylized via a text prompt.\lu{We can see a clear artifact in the middle row over the train (white block). This image may raise question.}}
\label{fig:outpaint_stylize}
\end{figure}

% A simple example of zero-shot editing is inpainting with text control as in~\Cref{fig:inpainting_dragon_goat}, but our model can do more by chaining multiple capabilities.
%
With the multi-task pretraining, \modelname{} exhibits task generalization that can be chained together to perform novel tasks.
%
We show the model can apply image-to-video animation followed by video-to-video stylization in \Cref{fig:animate_stylize}.
In the Appendix, \Cref{fig:outpaint_stylize} shows another example applying video-to-video outpainting, followed by editing them with additional video-to-video effects.
% In the supplementary materials, we include more examples.
%
At each stage, the quality of the output is sufficient to remain in-distribution (\ie teacher forcing) for the next stage without noticeable artifacts.
%
These capabilities can be attributed to our multimodal task design within an LLM transformer framework that allows for modeling multimodal content using a single transformer architecture over a unified vocabulary. 
%Our approach contrasts with others, such as diffusion models, which typically solve these tasks by adopting multiple individually tuned adapter models to control the diffusion process~\cite{esser2023structure,liew2023magicedit,guo2023animatediff}.

\vspace{-4mm}
\paragraph{Zero-shot video stylization.}
Stylization results are presented in \Cref{sec:stylization} where the structure and text are used as prefixes to guide the language model. Unlike other stylization methods that employ adapter modules such as cross-attention networks~\cite{zhang2023adding} or latent blending~\cite{meng2021sdedit}, our approach stylizes videos within an LLM as one of several generative tasks.
%We observe temporally coherent generations of objects in a video scene with dynamic, and meaningful motion (see \Cref{fig:long_video}).
% To predict the future frames, despite the model only being able to view up to a short temporal context, such as the first frame or the first second of video, the model is able to keep the motion, style, and identity of objects consistent across more than a total of 8 seconds of video output.


\vspace{-4mm}
\paragraph{3D structure, camera motion, visual styles.}
% Because our training spans videos, images, and text, we can prompt our model to demonstrate many aspects of understanding about the world including 3D structures, camera motions, and visual styles learned from these different sources.
Even though we do not add specific training data or losses to encourage 3D consistency, our model
can rotate around objects and predict reasonable visualizations of the backside of objects. Additionally, with only a small proportion of input videos and texts describing camera motion, our model can use short text prompts to apply a range of camera motions to image-to-video and text-to-video generations (see \Cref{fig:camera_movement}).

\subsection{Limitations}

Despite VideoPoet demonstrating highly competitive performance of LLMs relative to state-of-the-art models, certain limitations are still observed.
For example, the RGB frame reconstruction from compressed and quantized tokens place an upper bound on the generative model's visual fidelity. Second, the per-frame aesthetic biases in static scenes does not match the best baseline.
This difference is largely due to a choice of training data, where we focus our training on more natural aesthetics % and excluded some other sources common in other work.
and excluded some sources containing copyrighted images, such as LAION~\cite{schuhmann2022laion}, which is commonly used in other work.
% This shortfall is primarily due to the limited image training data, which is constrained to images for which we have rights or legal clearance to use, thereby excluding copyrighted images such as those found in the LAION dataset~\cite{schuhmann2022laion}.
Finally, small objects and fine-grained details, especially when coupled with significant motions, remains difficult within the token-based modeling.
%and so put an upper bound on the per-frame quality for a fixed amount of tokens.
% The video tokenizer is also trained on 8 fps video, which can lead to distored artifacts when reconstructing videos with fast motion.

% Per-frame quality is low due to the bad image\lu{complete this}.
%, which has been noted to be difficult for many existing video generation models~\cite{liu2023evalcrafter}. In addition, these controls can be added on top of a wide range of styles,  such as watercolor or oil paintings. 

% These stylization training sources are primarily observed in the text-image training data.
% The ability to generalize across and combine these different types of styles to produce large motions following text prompts underscores the strength of our model's understanding of objects in a temporal context.

